% 04-one-equation
\chapter{One equation is not enough}\label{ch:one-equation}

What can a single species do on its own? For one reaction--diffusion equation on
$(0,L)$ with Neumann conditions we prove two negative results: diffusion cannot
destabilize a constant steady state, and every nonconstant stationary solution is
linearly unstable. So one equation cannot produce linearly stable patterns. The tools
are elementary (integration by parts, Cauchy--Schwarz, cosine coefficients,
uniqueness for linear ODEs), and at the end we identify the structural reason behind
both results.

\section{Constant steady states}\label{sec:one-constant}

Let $d>0$ and let $f\colon\R\to\R$ be smooth. We consider
\begin{equation}\label{eq:scalar-rd}
\begin{aligned}
  u_t &= d\,u_{xx} + f(u), & x\in(0,L),\ t>0,\\
  u_x &= 0, & x\in\{0,L\},\ t>0,\\
  u(x,0) &= u_0(x), & x\in(0,L).
\end{aligned}
\end{equation}
If $f(\ub)=0$, the constant function $u\equiv\ub$ solves \eqref{eq:scalar-rd}; we call
it a \emph{constant steady state}. It is also a steady state of the \emph{kinetic
equation} $u'=f(u)$, asymptotically stable if $f'(\ub)<0$ and unstable if
$f'(\ub)>0$ (\cref{thm:linearized-stability}). We perturb, $u=\ub+w$, use
$f(\ub+w)=f'(\ub)w+\mathcal{O}(w^2)$ and drop the quadratic remainder, exactly as for
ODEs (the justification is discussed in \cref{ch:outlook}). This gives the
\emph{linearized problem}
\begin{equation}\label{eq:one-linearized}
  w_t = d\,w_{xx} + f'(\ub)\,w,\quad x\in(0,L),\ t>0,\qquad
  w_x = 0,\quad x\in\{0,L\},\ t>0 .
\end{equation}

\begin{definition}\label{def:linear-stability-scalar}
The constant steady state $\ub$ is \emph{linearly stable} if every solution of
\eqref{eq:one-linearized} (with smooth initial data satisfying the boundary
condition) satisfies $\norm{w(t)}\to0$ as $t\to\infty$. It is \emph{linearly
unstable} if \eqref{eq:one-linearized} has a solution with $\norm{w(t)}\to\infty$
as $t\to\infty$.
\end{definition}

\Cref{ch:turing} uses the same definition for systems (\cref{def:linear-stability-rd}).
We solve \eqref{eq:one-linearized} like the heat equation in \cref{ch:diffusion}. Let
$\hat w_k(t)$ be the cosine coefficients \eqref{eq:cosine-coefficients} of $w(\cdot,t)$.
Since $w_x=0$ at both ends, the coefficients of $w_{xx}$ are $-\lambda_k\hat w_k$
(\cref{lem:coefficient-uxx}). Differentiating under the integral sign for $t>0$, and
using that $\hat w_k$ is continuous at $t=0$, we get for $k=0,1,2,\dots$
\begin{equation}\label{eq:one-modes}
  \hat w_k'(t) = \big(f'(\ub)-d\lambda_k\big)\hat w_k(t),
  \qquad\text{hence}\qquad
  \hat w_k(t)=e^{\mu_kt}\hat w_k(0),\quad \mu_k = f'(\ub)-d\lambda_k .
\end{equation}
Since $\lambda_k\ge\lambda_0=0$, all growth rates satisfy $\mu_k\le\mu_0=f'(\ub)$.

\begin{proposition}\label{prop:scalar-constant}
Let $f(\ub)=0$. (a) If $f'(\ub)<0$, every solution of \eqref{eq:one-linearized}
satisfies $\norm{w(t)}\le e^{f'(\ub)t}\norm{w(0)}$ for $t\ge0$; in particular, $\ub$ is
linearly stable. (b) If $f'(\ub)>0$, then $\ub$ is linearly unstable.
\end{proposition}

\begin{proof}
(a) By Parseval's identity \eqref{eq:parseval}, \eqref{eq:one-modes} and
$\mu_k\le f'(\ub)$,
\begin{equation*}
  \norm{w(t)}^2 = Le^{2\mu_0t}\hat w_0(0)^2
     + \frac L2\sum_{k=1}^\infty e^{2\mu_kt}\hat w_k(0)^2
  \le e^{2f'(\ub)t}\norm{w(0)}^2 .
\end{equation*}
(b) The mode $k=0$, $w(x,t)=e^{f'(\ub)t}$, solves \eqref{eq:one-linearized}, and
$\norm{w(t)}=\sqrt L\,e^{f'(\ub)t}\to\infty$.
\end{proof}

We exclude the borderline case $f'(\ub)=0$. By analogy with \cref{def:ddi},
\emph{diffusion-driven instability (DDI)} would mean $f'(\ub)<0$ together with linear
instability for \eqref{eq:scalar-rd}.

\begin{corollary}\label{cor:no-ddi-scalar}
A single reaction--diffusion equation never exhibits diffusion-driven instability: if
$f'(\ub)<0$, then $\ub$ is linearly stable for every $d>0$ and every $L>0$.
\end{corollary}

\begin{proof}
This is \cref{prop:scalar-constant}(a).
\end{proof}

\begin{keyidea}
For one equation, the cosine mode $k$ grows at the rate $\mu_k=f'(\ub)-d\lambda_k$.
Diffusion only adds the damping $-d\lambda_k\le0$. The most dangerous perturbation is
the spatially constant one ($k=0$), which does not feel diffusion at all.
\end{keyidea}

\begin{example}\label{ex:one-kinetics}
For the logistic (Fisher--KPP) kinetics $f(u)=u(1-u)$, the state $0$ is linearly
unstable ($f'(0)=1$) and $1$ is linearly stable ($f'(1)=-1$). The bistable kinetics
$f(u)=u(1-u)(u-\theta)$, $0<\theta<1$, have $f'(0)=-\theta$,
$f'(\theta)=\theta(1-\theta)$, $f'(1)=\theta-1$: $0$ and $1$ are linearly stable,
$\theta$ is linearly unstable. Can a stationary solution close to $0$ in one part of the
interval and close to $1$ in another be linearly stable? No (\cref{sec:one-nonconstant}).
\end{example}

\section{A growth lemma}\label{sec:one-growth}

At a nonconstant stationary solution the linearized equation has an $x$-dependent
coefficient, and the cosine modes no longer decouple. We need a tool without Fourier
series. Let $a$ be a smooth function on $[0,L]$. We consider
\begin{equation}\label{eq:one-variable}
  w_t = d\,w_{xx} + a(x)\,w,\quad x\in(0,L),\ t>0,\qquad
  w_x=0,\quad x\in\{0,L\},\ t>0 ,
\end{equation}
and define $Q(\psi) = \int_0^L \big(a\,\psi^2 - d\,\psi'^2\big)\dd x$ for
$\psi\in C^1[0,L]$.

\begin{lemma}[Growth lemma]\label{lem:rayleigh-growth}
Let $w$ be a smooth solution of \eqref{eq:one-variable} (as in
\cref{rem:well-posedness}) such that $w(\cdot,0)$ is not identically zero, and put
$N(t)=\norm{w(t)}^2$.
\begin{enumerate}[label=(\alph*)]
  \item For $t>0$, $N'(t)=2Q(w(t))$ and $\frac{\mathrm d}{\mathrm dt}Q(w(t)) = 2\norm{w_t(t)}^2$.
  \item $N(t)>0$ for $t\ge0$, the \emph{Rayleigh quotient} $R(t)=Q(w(t))/N(t)$ is
        nondecreasing on $[0,\infty)$, and $\norm{w(t)}\ge e^{R(0)t}\norm{w(0)}$.
  \item $\norm{w(t)}\le e^{\alpha t}\norm{w(0)}$ for $t\ge0$, where $\alpha=\max_{[0,L]}a$.
\end{enumerate}
\end{lemma}

\begin{proof}
(a) For $t>0$ we differentiate under the integral sign, use the equation and integrate
by parts:
\begin{equation}\label{eq:one-N-prime}
  \frac12N' = \int_0^Lww_t\dd x = \int_0^Lw\,(dw_{xx}+aw)\dd x
  = d\big[ww_x\big]_0^L + \int_0^L\big(aw^2-dw_x^2\big)\dd x = Q(w),
\end{equation}
\begin{equation*}
  \frac{\mathrm d}{\mathrm dt}Q(w)
  = 2\int_0^L\big(aww_t - dw_xw_{xt}\big)\dd x
  = -2d\big[w_xw_t\big]_0^L + 2\int_0^L(dw_{xx}+aw)\,w_t\dd x
  = 2\int_0^Lw_t^2\dd x .
\end{equation*}
Both boundary terms vanish because $w_x=0$ at $x=0,L$. Moreover, $w$ and $w_x$ are
continuous on $[0,L]\times[0,\infty)$, hence uniformly continuous on
$[0,L]\times[0,1]$; so $N(t)\to N(0)$ and $Q(w(t))\to Q(w(0))$ as $t\downarrow0$.

(b) Let $T=\inf\{t>0: N(t)=0\}$ ($T=\infty$ if $N$ has no zero); $T>0$ since $N(0)>0$.
On $(0,T)$, part (a), the identity $Q(w)=\int_0^Lww_t\dd x$ from
\eqref{eq:one-N-prime} and the Cauchy--Schwarz inequality give
\begin{equation*}
  R' = \frac{N\,\frac{\mathrm d}{\mathrm dt}Q(w) - Q(w)\,N'}{N^2}
  = \frac{2}{N^2}\Big(\norm{w}^2\norm{w_t}^2
     - \Big(\int_0^Lww_t\dd x\Big)^2\Big)\ge0 .
\end{equation*}
As $R$ is continuous on $[0,T)$, it is nondecreasing there. Hence $(\log N)'=2R\ge2R(0)$
on $(0,T)$, and integrating (with $N$ continuous at $0$) gives $N(t)\ge N(0)e^{2R(0)t}$
on $[0,T)$. If $T$ were finite, letting $t\to T$ would give $N(T)>0$, contradicting
$N(T)=0$. So $T=\infty$, and taking square roots proves (b).

(c) $N'=2Q(w)\le2\int_0^Law^2\dd x\le2\alpha N$, so $e^{-2\alpha t}N(t)$ is
nonincreasing on $(0,\infty)$ and, by continuity, on $[0,\infty)$.
\end{proof}

By (b), to prove instability it suffices to find \emph{one} admissible initial datum
with $Q>0$. For constant $a\equiv f'(\ub)$, part (c) is the estimate of
\cref{prop:scalar-constant}(a), and $w(\cdot,0)\equiv1$ has $R(0)=f'(\ub)$: the
lemma reproves \cref{prop:scalar-constant} and \cref{cor:no-ddi-scalar} without
Fourier series.

\begin{figure}[t]
  \centering
  \includegraphics[width=\textwidth]{ch04-phase-plane}
  \caption{Stationary problem $U''+f(U)=0$, $U'(0)=U'(L)=0$, for
  $f(u)=u(1-u)(u-\frac12)$.
  (a) Level curves of $\frac12U'^2+F(U)$, with saddles (filled dots), center (open
  dot) and saddle connections (dashed); the highlighted orbit piece leaves the axis
  $U'=0$ at $x=0$ and returns at $x=L=8$. (b) The corresponding stationary solution;
  its linearization has a growing mode with rate $\approx0.081$ (numerically).}
  \label{fig:ch04-phase-plane}
\end{figure}

\section{Nonconstant stationary solutions are linearly unstable}\label{sec:one-nonconstant}

A \emph{stationary solution} of \eqref{eq:scalar-rd} is a function $U(x)$ with
\begin{equation}\label{eq:one-stationary}
  d\,U'' + f(U) = 0 \quad\text{in } (0,L),\qquad U'(0)=U'(L)=0 .
\end{equation}
We view \eqref{eq:one-stationary} as an ODE with $x$ as time. If $F'=f$, then
$\frac{\mathrm d}{\mathrm dx}\big(\frac d2U'^2+F(U)\big)=U'(dU''+f(U))=0$: the ODE is
conservative, and its orbits in the $(U,U')$ plane lie on level curves of
$\frac d2U'^2+F(U)$. A solution of \eqref{eq:one-stationary} is an orbit piece that
starts on the axis $U'=0$ at time $x=0$ and is on it again at time $x=L$.
\Cref{fig:ch04-phase-plane} shows this for the bistable kinetics with
$\theta=\frac12$, $d=1$. The orbit through $(U_0,0)$, $0<U_0<\frac12$, reaches
$(1-U_0,0)$ in half a period, which (numerically) increases from $2\pi$ for small
orbits to $\infty$ near the saddle connections. So for every $L>2\pi$ there is an
increasing nonconstant stationary solution.

Linearizing \eqref{eq:scalar-rd} at a stationary solution $U$, $u=U+w$, gives
\begin{equation}\label{eq:one-linearized-U}
  w_t = d\,w_{xx} + f'(U(x))\,w,\quad x\in(0,L),\ t>0,\qquad
  w_x=0,\quad x\in\{0,L\},\ t>0 ,
\end{equation}
that is, \eqref{eq:one-variable} with $a=f'(U)$; $Q$ denotes the corresponding form.
We call $U$ linearly (un)stable as in \cref{def:linear-stability-scalar}, with
\eqref{eq:one-linearized-U} in place of \eqref{eq:one-linearized}. Note that $U$ is
smooth, since $U''=-f(U)/d$.

\begin{theorem}\label{thm:nonconstant-unstable}
Every nonconstant solution $U$ of \eqref{eq:one-stationary} is linearly unstable.
More precisely, \eqref{eq:one-linearized-U} has a solution with
$\norm{w(t)}\ge c\,e^{\mu t}$ for all $t\ge0$, where $c,\mu>0$ are constants.
\end{theorem}

\begin{proof}
\emph{Step 1: reduction.} Let $w_0$ be smooth on $[0,L]$ with $w_0'(0)=w_0'(L)=0$ and
$Q(w_0)>0$, and let $w$ solve \eqref{eq:one-linearized-U} with $w(\cdot,0)=w_0$
(\cref{rem:well-posedness}). By \cref{lem:rayleigh-growth}(b),
$\norm{w(t)}\ge e^{\mu t}\norm{w_0}$ with $\mu=R(0)=Q(w_0)/\norm{w_0}^2>0$. It remains
to find such a $w_0$.

\emph{Step 2: the derivative of $U$.} Put $z=U'$. Differentiating $dU''+f(U)=0$ gives
$dz''+az=0$, and the Neumann condition for $U$ says $z(0)=z(L)=0$. Hence
\begin{equation*}
  Q(z) = \int_0^L\big(az^2-dz'^2\big)\dd x
  = \int_0^Lz\,(dz''+az)\dd x - d\big[zz'\big]_0^L = 0 .
\end{equation*}
But $z$ does not qualify: $Q(z)=0$, and $z'(0)=-f(U(0))/d$ need not vanish.

\emph{Step 3: making $Q$ positive.} First, $z'(0)\neq0$: otherwise $z$ would solve
the linear ODE $dz''+az=0$ with $z(0)=z'(0)=0$, so $z\equiv0$ by uniqueness, and $U$
would be constant. For smooth $\eta$ and $\eps\in\R$, expanding and using $Q(z)=0$
gives $Q(z+\eps\eta)=2\eps B+\eps^2Q(\eta)$, where, by parts and since $dz''+az=0$,
\begin{equation*}
  B=\int_0^L\big(az\eta-dz'\eta'\big)\dd x
  = \int_0^L\eta\,(dz''+az)\dd x - d\big[z'\eta\big]_0^L
  = d\big(z'(0)\eta(0)-z'(L)\eta(L)\big).
\end{equation*}
For $\eta(x)=z'(0)(1-x/L)$ we get $Q(z+\eps\eta)=2\eps d\,z'(0)^2+\eps^2Q(\eta)>0$ for
small $\eps>0$. Fix such an $\eps$; then $\psi=z+\eps\eta$ is smooth with $Q(\psi)>0$,
but need not satisfy the Neumann condition.

\emph{Step 4: the Neumann condition.} For $0<\tau<L/2$ let
$g_\tau\colon[0,L]\to[0,L]$ be smooth and nondecreasing with $g_\tau(0)=0$,
$g_\tau(L)=L$, $g_\tau'(0)=g_\tau'(L)=0$, $g_\tau(x)=x$ on $[\tau,L-\tau]$
and $0\le g_\tau'\le2$ (\cref{exer:one-g-delta}). Then $\psi_\tau=\psi\circ g_\tau$
is smooth, $\psi_\tau'(0)=\psi'(0)g_\tau'(0)=0$, likewise $\psi_\tau'(L)=0$,
and $\psi_\tau=\psi$ outside $I_\tau=[0,\tau]\cup[L-\tau,L]$. With
$M_0=\max\abs{\psi}$, $M_1=\max\abs{\psi'}$, $C_a=\max\abs{a}$ we have
$\abs{\psi_\tau}\le M_0$, $\abs{\psi_\tau'}\le2M_1$, so
\begin{equation*}
  \abs{Q(\psi_\tau)-Q(\psi)}
  \le \int_{I_\tau}\big(\abs{a}\psi_\tau^2+d\psi_\tau'^2
      +\abs{a}\psi^2+d\psi'^2\big)\dd x
  \le 2\tau\big(2C_aM_0^2+5dM_1^2\big).
\end{equation*}
Hence $Q(\psi_\tau)>0$ for small $\tau$, and $w_0=\psi_\tau$ does the job.
\end{proof}

\begin{keyidea}
In one space dimension with Neumann conditions, every nonconstant stationary solution
of a single reaction--diffusion equation is linearly unstable. The only candidates for
linearly stable stationary states are constants.
\end{keyidea}

\section{Remarks}\label{sec:one-remarks}

\begin{remark}[An energy]\label{rem:lyapunov}
Let $F'=f$ and $E[u]=\int_0^L\big(\frac d2u_x^2-F(u)\big)\dd x$. If $u$ solves
\eqref{eq:scalar-rd}, then for $t>0$, integrating by parts and using $u_x=0$ at the ends,
\begin{equation*}
  \frac{\mathrm d}{\mathrm dt}E[u(t)] = \int_0^L\big(d\,u_xu_{xt}-f(u)u_t\big)\dd x
  = d\big[u_xu_t\big]_0^L - \int_0^L\big(du_{xx}+f(u)\big)u_t\dd x = -\norm{u_t}^2 ,
\end{equation*}
and $E[u(t)]\to E[u_0]$ as $t\downarrow0$, since $u$, $u_x$ are continuous up to $t=0$.
So $E$ is a \emph{Lyapunov functional}: it decreases along solutions and is constant
only along stationary ones (hence time-periodic solutions are stationary). For
stationary $U$ and smooth $\eta$, Taylor's formula and the same integration by parts
give $E[U+\eps\eta]=E[U]-\frac{\eps^2}2Q(\eta)+\mathcal{O}(\eps^3)$ with $a=f'(U)$:
stationary solutions are critical points of $E$, and $-Q$ is the second variation.
If $U$ is nonconstant and $w_0$ is as in the proof of
\cref{thm:nonconstant-unstable}, then for small $\eps>0$ the datum
$u_0=U+\eps w_0$ gives $E[u(t)]\le E[u_0]<E[U]$ for all $t\ge0$, so $u(t)$ cannot
converge to $U$ in any sense in which $E$ is continuous, e.g.\ in the distance
$\max\abs{u-U}+\max\abs{u_x-U'}$. As $u_0$ is arbitrarily close to $U$ in this
distance, $U$ is \emph{not asymptotically stable} for the nonlinear equation. This
is weaker than instability (solutions starting near $U$ might stay near it); under
suitable assumptions on $f$, nonlinear instability holds as well
\cite{CastenHolland1978,Chafee1975}.
\end{remark}

\begin{remark}[Higher dimensions]\label{rem:higher-dim}
One dimension was used in Step 2: $z=U'$ solves the linearized equation \emph{and}
vanishes at the boundary. On $\Omega\subset\R^n$ the derivatives $\partial_jU$ solve
the linearized equation but need not vanish on $\partial\Omega$. On convex domains a
careful treatment of the boundary terms still gives instability
\cite{CastenHolland1978,Matano1979}. On a dumbbell (two regions joined by a thin
channel), however, bistable kinetics can have stable stationary solutions close to
$0$ in one bulb and to $1$ in the other \cite{Matano1979}. In one dimension the
long-time behavior is well understood \cite{Chafee1975}.
\end{remark}

\begin{remark}[Traveling waves]\label{rem:travelling-waves}
On the whole line, $u_t=du_{xx}+u(1-u)$ has \emph{traveling fronts}
$u=\Phi(x-ct)$, $0<\Phi<1$, $\Phi(-\infty)=1$, $\Phi(+\infty)=0$, exactly for
$c\ge2\sqrt d$ \cite{Murray2002}. On a bounded interval with Neumann conditions a
front is only transient: it crosses the interval, and then the solution approaches
the constant $1$ (for nonnegative data, not identically zero).
\end{remark}

\section{Why two species?}\label{sec:one-why-two}

\Cref{thm:nonconstant-unstable} rests on the symmetry
$\int_0^L(d\psi''+a\psi)\eta\dd x=\int_0^L\psi(d\eta''+a\eta)\dd x$ for smooth
$\psi,\eta$ with $\psi'=\eta'=0$ at $x=0,L$ (integrate by parts twice), which gave
$\frac{\mathrm d}{\mathrm dt}Q(w)=2\norm{w_t}^2$; \cref{cor:no-ddi-scalar} rests on the
fact that a scalar reaction term satisfies $aw^2\le(\max a)w^2$. For two species both
properties hold if $A$ is symmetric, $f_v=g_u$: then $\bw\mapsto D\bw_{xx}+A\bw$ from
\eqref{eq:linearized-system} is symmetric in the same sense, and diffusion is harmless:

\begin{proposition}\label{prop:symmetric-no-ddi}
Let $A$ be symmetric with two negative eigenvalues, the larger one $\mu_A$, and
$D=\diag(d_u,d_v)$, $d_u,d_v>0$. Then $\mathbf y\cdot(A-\lambda D)\mathbf y\le
\mu_A\abs{\mathbf y}^2$ for all $\lambda\ge0$, $\mathbf y\in\R^2$ (every $A-\lambda D$
is negative definite), and every solution of \eqref{eq:linearized-system} satisfies
$\norm{\bw(t)}\le e^{\mu_At}\norm{\bw(0)}$. So there is no diffusion-driven instability.
\end{proposition}

\begin{proof}
Expanding $\mathbf y$ in an orthonormal eigenbasis of $A$ gives
$\mathbf y\cdot A\mathbf y\le\mu_A\abs{\mathbf y}^2$, and
$\mathbf y\cdot D\mathbf y=d_uy_1^2+d_vy_2^2\ge0$. As in \cref{lem:rayleigh-growth}(c),
integrating by parts (the boundary term vanishes),
$\frac{\mathrm d}{\mathrm dt}\norm{\bw}^2=-2\int_0^L\bw_x\cdot D\bw_x\dd x
+2\int_0^L\bw\cdot A\bw\dd x\le2\mu_A\norm{\bw}^2$.
\end{proof}

\begin{remark}[Why two species]\label{rem:why-two}
One equation admits neither diffusion-driven instability (\cref{cor:no-ddi-scalar})
nor linearly stable nonconstant stationary solutions (\cref{thm:nonconstant-unstable}),
and symmetric two-species systems admit no diffusion-driven instability
(\cref{prop:symmetric-no-ddi}). A Turing instability needs two species and a
\emph{non-symmetric} interaction, $f_v\neq g_u$. \Cref{ch:turing} shows more:
$f_vg_u<0$ is necessary (\cref{cor:turing-consequences}), so one species promotes the
other while being suppressed by it. The last proof only used
$\mathbf y\cdot A\mathbf y\le\mu_A\abs{\mathbf y}^2$. For non-symmetric $A$, stable
eigenvalues do not imply $\mathbf y\cdot A\mathbf y<0$; an instability needs
$\mathbf y\cdot A\mathbf y>0$ for some $\mathbf y$, e.g.\ $f_u>0$ (self-enhancement).
\end{remark}

\section{Exercises}\label{sec:one-exercises}

\begin{exercise}\label{exer:one-fisher}
Let $f(u)=u(1-u)$. Show that at $\ub=0$ exactly the modes $0\le k<L/(\pi\sqrt d)$ grow,
the fastest being $k=0$; list them with their rates for $d=1$, $L=10$. Show that at
$\ub=1$ every mode decays at least like $e^{-t}$. Compare with
\cref{ex:schnakenberg-turing}, where $k=6$ grows fastest.
\end{exercise}

\begin{exercise}\label{exer:one-g-delta}
Let $\rho$ be smooth and nondecreasing with $\rho=0$ on $(-\infty,0]$, $\rho=1$ on
$[1,\infty)$ and $\rho(s)+\rho(1-s)=1$ (e.g.\
$\rho(s)=e^{-1/s}/(e^{-1/s}+e^{-1/(1-s)})$ on $(0,1)$). Let
$H(s)=\rho(2s)+\rho(4s-2)\rho(4-4s)$ for $0\le s\le1$, $H(s)=1$ for $s>1$, and
$G(s)=\int_0^sH$. Show that $0\le H\le2$, $H(0)=0$, $\int_0^1H=1$, and that
$g_\tau(x)=\tau G(x/\tau)$ for $x\le L/2$, $g_\tau(x)=L-\tau G((L-x)/\tau)$
for $x>L/2$, has all properties required in Step 4.
\end{exercise}

\begin{exercise}[$\ast$]\label{exer:one-minimal-length}
Let $f(u)=u-u^3$, $F(u)=\frac12u^2-\frac14u^4$. (a) Show that every nonconstant
solution of \eqref{eq:one-stationary} lies on a closed orbit through $(\pm U_m,0)$,
$0<U_m<1$, and $L=mT(U_m)/2$ with $m\in\N$, where $T(U_m)$ is the period.
(b) Show $T(U_m)=\sqrt{2d}\int_{-U_m}^{U_m}(F(U_m)-F(U))^{-1/2}\dd U>2\pi\sqrt d$, using
$F(U_m)-F(U)<\frac12(U_m^2-U^2)$. (c) Assuming that $T$ is continuous,
$T(U_m)\to2\pi\sqrt d$ as $U_m\to0$ and $T(U_m)\to\infty$ as $U_m\to1$, show that a
nonconstant stationary solution exists if and only if $L>\pi\sqrt d$. Which mode of
$\ub=0$ starts to grow there?
\end{exercise}

\begin{exercise}[$\ast$]\label{exer:one-dirichlet}
Consider $u_t=u_{xx}+u(1-u)$ on $(0,L)$ with $u(0,t)=u(L,t)=0$. (a) As in
\cref{exer:one-minimal-length}, show that for $L>\pi$ there is a stationary solution
$U$ with $U(0)=U(L)=0$, $0<U<1$ in $(0,L)$. (b) Show that $N'=2Q(w)$, with $a=1-2U$,
$d=1$, still holds for the linearization $w_t=w_{xx}+(1-2U)w$, $w(0,t)=w(L,t)=0$.
(c) Let $w\not\equiv0$ be smooth with $w(0)=w(L)=0$. Show that $U'(0),U'(L)\neq0$
(uniqueness for the ODE) and deduce that $\psi=w/U$ is smooth on $[0,L]$ (near $0$,
$\psi(x)=\int_0^1w'(sx)\dd s/\int_0^1U'(sx)\dd s$). Show
$\int_0^L(w'^2-(1-U)w^2)\dd x=\int_0^LU^2\psi'^2\dd x$ and deduce
$Q(w)\le-\int_0^LUw^2\dd x<0$: no perturbation grows. (In fact $U$ is linearly
stable.) (d) Which step of the proof of \cref{thm:nonconstant-unstable} fails?
\end{exercise}
